% ============================================================================
%  CAPÍTULO III: DESARROLLO DEL TRABAJO (parte 1)
% ============================================================================

\section{Plan de migración}

\subsection{Estrategia seleccionada}

La selección de estrategia dependió de tres variables: el volumen (0.28 MB y
1\,469 filas), la ventana de mantenimiento disponible y la criticidad del
negocio. Se evaluaron tres alternativas. La primera, migración \emph{big bang}
con corte controlado, consiste en congelar escrituras en el origen, cargar los
datos a un schema de staging (\texttt{stg}) con pgloader, validar con
\emph{precheck}, aplicar en una sola transacción con rollback garantizado, y
cambiar la cadena de conexión. La segunda, replicación continua con
\emph{Change Data Capture} (CDC), mantiene ambos motores sincronizados y
reduce el tiempo de inactividad, pero requiere herramientas adicionales,
monitoreo permanente y un esfuerzo de ingeniería desproporcionado para el
volumen actual. La tercera, migración por lotes con doble escritura, exige
modificar la aplicación y convivir con dos fuentes de verdad durante semanas.

La Tabla~\ref{tab:estrategias} compara las alternativas. Para este proyecto se
eligió el corte controlado con staging porque el volumen permite copiar y
verificar en segundos, el respaldo y la restauración están probados, y el
rollback es transaccional: si el \emph{precheck} o el \emph{postcheck} falla,
se conserva el origen y el destino no sufre modificaciones parciales. La
ventana de mantenimiento estimada es de 30 a 60 minutos, muy por debajo de las
4 horas que justificarían CDC.

\begin{table}[H]
  \centering
  \caption{Comparación de estrategias de migración}
  \label{tab:estrategias}
  \contenidoConFuente{%
    \begin{tablaAPA}
      \begin{tabular}{@{}L{3.2cm}L{3.4cm}L{3.4cm}L{4.2cm}@{}}
        \toprule
        \textbf{Estrategia} & \textbf{Ventaja} & \textbf{Desventaja} & \textbf{Decisión} \\
        \midrule
        \emph{Big bang} controlado con staging & Simple, verificable, rollback transaccional; pgloader + precheck/postcheck. & Requiere ventana de inactividad y congelamiento de escrituras. & \textbf{Seleccionada}: el volumen y la criticidad lo permiten. \\
        CDC con replicación & Inactividad mínima; útil para volúmenes grandes. & Complejidad operativa y de monitoreo; herramientas adicionales. & Descartada para esta fase. \\
        Doble escritura & Transición gradual. & Cambios en la aplicación; riesgo de divergencia. & Descartada: exige modificar el código. \\
        \bottomrule
      \end{tabular}
    \end{tablaAPA}%
  }{Elaboración propia.}
\end{table}

\subsection{Fases y cronograma}

El plan se organizó en diez fases. La Tabla~\ref{tab:fases} indica el
objetivo, las actividades principales, el entregable y la duración estimada.
Las fases 1 a 10 se ejecutaron en el laboratorio; las fases de validación se
repitieron hasta obtener 75 pruebas en \texttt{PASS}.

\begin{table}[H]
  \centering
  \caption{Fases del plan de migración}
  \label{tab:fases}
  \contenidoConFuente{%
    \begin{tablaAPA}
      \begin{tabular}{@{}C{1.0cm}L{2.9cm}L{6.2cm}L{3.4cm}C{1.3cm}@{}}
        \toprule
        \textbf{Fase} & \textbf{Nombre} & \textbf{Actividades} & \textbf{Entregable} & \textbf{Duración} \\
        \midrule
        1 & Congelamiento & Detener escrituras en MySQL; validar conectividad origen--destino. & Entorno \emph{healthy} congelado. & 5 min \\
        2 & Respaldo origen & \texttt{mysqldump --single-transaction} del origen. & \texttt{backups/origin\_*.sql.gz}. & 10 min \\
        3 & Staging & \texttt{00\_prepare\_staging.sql} crea schema \texttt{stg}. & Schema \texttt{stg} listo. & 5 min \\
        4 & Carga pgloader & pgloader carga 8 tablas a \texttt{stg} (sin índices, FK, secuencias). & Datos en \texttt{stg}. & 10 min \\
        5 & Conteo staging & Conteos origen vs \texttt{stg}. & Conteos coincidentes. & 5 min \\
        6 & Precheck & \texttt{10\_precheck.sql} valida nulos, ENUM, huérfanos, rangos, fechas cero. Aborta si falla. & Validación superada. & 10 min \\
        7 & Carga transaccional & \texttt{20\_load.sql} + \texttt{30\_postcheck.sql} + \texttt{35\_sequences.sql} en una transacción (\texttt{ROLLBACK} si falla). & Datos en \texttt{public} verificados. & 10 min \\
        8 & Estadísticas & \texttt{40\_analyze.sql} como \texttt{postgres} por socket local. & Planificador actualizado. & 5 min \\
        9 & Limpieza & \texttt{50\_cleanup.sql} elimina schema \texttt{stg}. & Entorno limpio. & 5 min \\
        10 & Respaldo destino y restauración & \texttt{pg\_dump -Fc} + AES-256-CBC; prueba de restauración con 55 socios. & Respaldo cifrado verificado. & 15 min \\
        \bottomrule
      \end{tabular}
    \end{tablaAPA}%
  }{Elaboración propia.}
\end{table}

\subsection{Plan de contingencia y reversión}

El plan de contingencia se diseñó antes de ejecutar la migración. La lógica es
simple: el origen es la fuente de verdad hasta que el destino pasa todas las
pruebas; si algo falla, se conserva el origen y se restaura el destino. La
Tabla~\ref{tab:contingencia} define los disparadores, las acciones y los
responsables. El objetivo de tiempo de recuperación (RTO) es de 30 minutos y
el objetivo de punto de recuperación (RPO) es de 0 minutos para los datos
migrados, porque el respaldo se toma inmediatamente antes del corte.

\begin{table}[H]
  \centering
  \caption{Plan de contingencia y reversión}
  \label{tab:contingencia}
  \contenidoConFuente{%
    \begin{tablaAPA}
      \begin{tabular}{@{}L{3.5cm}L{6.6cm}L{4.6cm}@{}}
        \toprule
        \textbf{Disparador} & \textbf{Acción} & \textbf{Responsable} \\
        \midrule
        Conteos origen vs \texttt{stg} no coinciden & Detener, revisar la carga pgloader, restaurar el destino desde el respaldo y repetir. & DBA / ETL \\
        \texttt{PRECHECK FALLIDO} & Abortar antes de tocar \texttt{public}; corregir la validación fallida (nulos, ENUM, huérfanos, rangos, fechas cero) y repetir desde la fase de carga. & DBA \\
        Error en transacción única & \texttt{ROLLBACK} automático; el destino no sufre modificaciones parciales. Revisar el error en \texttt{20\_load.sql} o \texttt{30\_postcheck.sql} y repetir. & DBA \\
        Clave foránea violada & Identificar el huérfano, corregir la causa en el origen o en la transformación, repetir la carga. & DBA \\
        Falla de conexión TLS & Revisar certificados, \texttt{pg\_hba} y \texttt{sslrootcert}; no continuar sin cifrado. & Seguridad \\
        Error de secuencia & Ejecutar \texttt{setval} con el máximo ID y probar un \texttt{INSERT} nuevo. & DBA \\
        Pérdida de datos & Restaurar el respaldo cifrado, validar 55 socios y 1\,469 filas, repetir la migración. & DBA \\
        Incidente de seguridad & Aislar el contenedor, revocar credenciales, revisar logs y notificar. & Seguridad / TI \\
        \bottomrule
      \end{tabular}
    \end{tablaAPA}%
  }{Elaboración propia.}
\end{table}

Los criterios de aceptación fueron: (a) 100\,\% de las filas migradas y
verificadas; (b) cero huérfanos en las 9 relaciones; (c) 75 pruebas en
\texttt{PASS}; (d) conexiones sin SSL rechazadas; (e) \texttt{auditor\_ro} sin
acceso a tablas sensibles; (f) respaldo cifrado restaurado con 55 socios; y
(g) entorno levantado desde cero con un solo comando. Si cualquiera de estos
criterios falla, la migración no se considera aceptada.

\section{Diseño del esquema destino}

\subsection{Principios de diseño}

El esquema destino conserva los nombres de las tablas y columnas del origen
para reducir el impacto en la aplicación, pero cambia los tipos y añade
restricciones. Se aplicaron cinco principios. Primero, usar tipos nativos de
PostgreSQL en lugar de emular los de MySQL. Segundo, declarar la integridad
referencial en el motor. Tercero, mantener la precisión decimal de los montos.
Cuarto, indexar las columnas usadas en reportes y claves foráneas. Quinto,
exponer los datos sensibles solo a través de vistas controladas.

La Figura~\ref{fig:modelo-datos} muestra el modelo relacional simplificado.
Las relaciones principales son: un socio tiene muchas cuentas y muchas
solicitudes; una solicitud aprobada genera un préstamo; un préstamo tiene
muchas cuotas; y las transacciones referencian un socio, una cuenta o un
préstamo según el tipo de movimiento.

\figuraAPA{fig:modelo-datos}{Modelo relacional simplificado del destino}{%
  \begin{tikzpicture}[
      ent/.style={draw, rounded corners, minimum width=2.7cm,
                  minimum height=0.85cm, align=center, font=\scriptsize,
                  fill=blue!4},
      linea/.style={-Stealth, thick, black!70},
      etiqueta/.style={font=\tiny, fill=white, inner sep=1pt}]
    \node[ent] (suc) at (0,0) {\texttt{sucursales}};
    \node[ent] (soc) at (3.4,0) {\texttt{socios}};
    \node[ent] (cta) at (6.8,0) {\texttt{cuentas\_ahorro}};
    \node[ent] (tx)  at (10.2,0) {\texttt{transacciones}};
    \node[ent] (tip) at (0,-2.6) {\texttt{tipos\_credito}};
    \node[ent] (sol) at (3.4,-2.6) {\texttt{solicitudes\_credito}};
    \node[ent] (pre) at (6.8,-2.6) {\texttt{prestamos}};
    \node[ent] (cuo) at (10.2,-2.6) {\texttt{cuotas\_amortizacion}};

    \draw[linea] (soc) -- node[etiqueta,above]{1:N} (cta);
    \draw[linea] (soc) -- node[etiqueta,right]{1:N} (sol);
    \draw[linea] (tip) -- node[etiqueta,above]{1:N} (sol);
    \draw[linea] (sol) -- node[etiqueta,above]{1:1} (pre);
    \draw[linea] (pre) -- node[etiqueta,above]{1:N} (cuo);
    \draw[linea] (cta) -- node[etiqueta,above]{1:N} (tx);
    \draw[linea] (pre) to[bend left=18] node[etiqueta,right]{1:N} (tx);
    \draw[linea, dashed] (soc) to[bend left=12] node[etiqueta,left]{1:N} (pre);
  \end{tikzpicture}%
}{Elaboración propia con base en \texttt{src/target/01\_schema\_target.sql}.}

\subsection{Tablas, claves y restricciones}

La Tabla~\ref{tab:tablas-destino} describe las 8 tablas, su clave primaria,
sus claves foráneas y las restricciones adicionales. Todas las claves
primarias usan identidades \texttt{BIGINT}. Las claves foráneas tienen
políticas explícitas: \texttt{ON DELETE CASCADE} en cuotas (si se elimina un
préstamo, sus cuotas no tienen sentido) y \texttt{ON DELETE SET NULL} en
transacciones (el historial debe conservarse aunque se elimine la referencia).

\begin{table}[H]
  \centering
  \caption{Tablas y restricciones del esquema destino}
  \label{tab:tablas-destino}
  \contenidoConFuente{%
    \begin{tablaAPA}
      \begin{tabular}{@{}L{3.0cm}L{3.2cm}L{4.0cm}L{4.2cm}@{}}
        \toprule
        \textbf{Tabla} & \textbf{Clave primaria} & \textbf{Claves foráneas} & \textbf{Restricciones} \\
        \midrule
        \texttt{sucursales} & \texttt{id} & --- & \texttt{estatus BOOLEAN}. \\
        \texttt{socios} & \texttt{id} & --- & \texttt{UNIQUE(numero\_identificacion)}; \texttt{CHECK(genero)}; \texttt{CHECK(estado)}. \\
        \texttt{cuentas\_ahorro} & \texttt{id} & \texttt{socio\_id} $\rightarrow$ \texttt{socios} & \texttt{UNIQUE(num\_cuenta)}; saldo y tasa no negativos. \\
        \texttt{tipos\_credito} & \texttt{id} & --- & \texttt{CHECK(estado)}; monto y plazo positivos. \\
        \texttt{solicitudes\_credito} & \texttt{id} & \texttt{socio\_id}, \texttt{tipo\_credito\_id} & \texttt{CHECK(status)}; monto y plazo positivos. \\
        \texttt{prestamos} & \texttt{id} & \texttt{solicitud\_id}, \texttt{socio\_id} & \texttt{UNIQUE(solicitud\_id)}; \texttt{CHECK(estado)}. \\
        \texttt{cuotas\_amortizacion} & \texttt{id} & \texttt{prestamo\_id} & \texttt{UNIQUE(prestamo\_id, numero\_cuota)}; \texttt{CHECK(estatus)}. \\
        \texttt{transacciones} & \texttt{id} & \texttt{socio\_id}, \texttt{cuenta\_ahorro\_id}, \texttt{prestamo\_id} & \texttt{CHECK(tipo)}; monto positivo. \\
        \bottomrule
      \end{tabular}
    \end{tablaAPA}%
  }{Elaboración propia con datos de \texttt{src/target/01\_schema\_target.sql}.}
\end{table}

El fragmento de DDL de la Tabla~\ref{tab:tablas-destino} se implementó en
\texttt{src/target/01\_schema\_target.sql}. Se reproducen las definiciones de
\texttt{socios}, \texttt{prestamos} y \texttt{transacciones} para mostrar el
estilo.

\begin{lstlisting}[language=SQL,caption={Definiciones representativas del esquema destino},label={lst:ddl-destino}]
CREATE TABLE socios (
  id BIGINT GENERATED ALWAYS AS IDENTITY PRIMARY KEY,
  nombre_completo VARCHAR(150) NOT NULL,
  numero_identificacion VARCHAR(20) UNIQUE,
  fecha_nacimiento DATE,
  genero VARCHAR(1) NOT NULL DEFAULT 'M'
    CHECK (genero IN ('M', 'F', 'O')),
  estado VARCHAR(10) NOT NULL DEFAULT 'activo'
    CHECK (estado IN ('activo', 'inactivo', 'suspendido')),
  fecha_ingreso DATE NOT NULL,
  telefono VARCHAR(20),
  email VARCHAR(100),
  direccion_residencia TEXT,
  fecha_ultima_actualizacion TIMESTAMPTZ NOT NULL DEFAULT now(),
  estatus BOOLEAN NOT NULL DEFAULT TRUE
);

CREATE TABLE prestamos (
  id BIGINT GENERATED ALWAYS AS IDENTITY PRIMARY KEY,
  solicitud_id BIGINT NOT NULL UNIQUE REFERENCES solicitudes_credito(id),
  socio_id BIGINT NOT NULL REFERENCES socios(id),
  monto_prestado NUMERIC(12,2) NOT NULL CHECK (monto_prestado > 0),
  tasa_interes_anual NUMERIC(5,2) NOT NULL CHECK (tasa_interes_anual >= 0),
  fecha_desembolso TIMESTAMPTZ NOT NULL,
  fecha_fin_plazo DATE NOT NULL,
  estado VARCHAR(10) NOT NULL DEFAULT 'activo'
    CHECK (estado IN ('activo', 'pagado', 'moroso', 'cancelado')),
  total_cuotas BIGINT NOT NULL CHECK (total_cuotas > 0),
  pagado_total BOOLEAN NOT NULL DEFAULT FALSE
);

CREATE TABLE transacciones (
  id BIGINT GENERATED ALWAYS AS IDENTITY PRIMARY KEY,
  tipo VARCHAR(15) NOT NULL
    CHECK (tipo IN ('deposito', 'retiro', 'pago_cuota', 'desembolso', 'transferencia')),
  socio_id BIGINT REFERENCES socios(id) ON DELETE SET NULL,
  cuenta_ahorro_id BIGINT REFERENCES cuentas_ahorro(id) ON DELETE SET NULL,
  prestamo_id BIGINT REFERENCES prestamos(id) ON DELETE SET NULL,
  monto NUMERIC(12,2) NOT NULL CHECK (monto > 0),
  fecha_transaccion TIMESTAMPTZ NOT NULL,
  descripcion VARCHAR(255)
);
\end{lstlisting}

\subsection{Índices}

Se crearon 10 índices no únicos además de las claves primarias y las
restricciones únicas. La selección se basó en las consultas de los reportes y
en las columnas de join. La Tabla~\ref{tab:indices-destino} justifica cada
índice. No se indexó todo: cada índice adicional consume espacio y ralentiza
las escrituras. Se priorizaron las columnas con mayor selectividad y uso.

\begin{table}[H]
  \centering
  \caption{Índices del esquema destino y su justificación}
  \label{tab:indices-destino}
  \contenidoConFuente{%
    \begin{tablaAPA}
      \begin{tabular}{@{}L{4.2cm}L{3.2cm}L{7.0cm}@{}}
        \toprule
        \textbf{Índice} & \textbf{Tabla} & \textbf{Uso previsto} \\
        \midrule
        \texttt{idx\_socio\_estado} & \texttt{socios} & Filtrar socios activos, inactivos o suspendidos. \\
        \texttt{idx\_cuenta\_ahorro\_socio} & \texttt{cuentas\_ahorro} & Reporte de saldos por socio. \\
        \texttt{idx\_solicitud\_tipo} & \texttt{solicitudes\_credito} & Análisis por producto de crédito. \\
        \texttt{idx\_solicitud\_socio} & \texttt{solicitudes\_credito} & Historial de solicitudes por socio. \\
        \texttt{idx\_transaccion\_fecha} & \texttt{transacciones} & Reportes por rango de fechas. \\
        \texttt{idx\_transaccion\_socio} & \texttt{transacciones} & Movimientos por socio. \\
        \texttt{idx\_prestamo\_estado} & \texttt{prestamos} & Cartera por estado (incluye morosos). \\
        \texttt{idx\_prestamo\_socio} & \texttt{prestamos} & Préstamos por socio. \\
        \texttt{idx\_cuota\_prestamo} & \texttt{cuotas\_amortizacion} & Plan de pagos por préstamo. \\
        \texttt{idx\_cuota\_estatus} & \texttt{cuotas\_amortizacion} & Cuotas pendientes, pagadas o atrasadas. \\
        \bottomrule
      \end{tabular}
    \end{tablaAPA}%
  }{Elaboración propia con datos de \texttt{src/target/01\_schema\_target.sql}.}
\end{table}

\subsection{Vistas}

Se recrearon las dos vistas del origen y se añadió la vista de auditoría. La
vista \texttt{vw\_cartera\_mora} conserva la lógica original: selecciona
préstamos en estado \texttt{moroso}, cuenta sus cuotas y suma las atrasadas.
La vista \texttt{vw\_resumen\_saldos\_socio} agrupa cuentas y préstamos por
socio. La vista \texttt{vw\_socios\_audit} es nueva y aplica enmascaramiento.
La Tabla~\ref{tab:vistas-destino} resume las tres.

\begin{table}[H]
  \centering
  \caption{Vistas del esquema destino}
  \label{tab:vistas-destino}
  \contenidoConFuente{%
    \begin{tablaAPA}
      \begin{tabular}{@{}L{4.4cm}L{2.4cm}L{7.4cm}@{}}
        \toprule
        \textbf{Vista} & \textbf{Tipo} & \textbf{Propósito} \\
        \midrule
        \texttt{vw\_cartera\_mora} & Reporte & Préstamos morosos, cuotas atrasadas y monto en mora. \\
        \texttt{vw\_resumen\_saldos\_socio} & Reporte & Saldos activos y cerrados, cantidad de cuentas y préstamos. \\
        \texttt{vw\_socios\_audit} & Auditoría & Datos de socio con CI/NIT, teléfono, correo y cuentas enmascarados. \\
        \bottomrule
      \end{tabular}
    \end{tablaAPA}%
  }{Elaboración propia con datos de \texttt{src/target/01\_schema\_target.sql} y \texttt{02\_hardening\_roles.sql}.}
\end{table}

\section{Implementación del ETL}

\subsection{Arquitectura del proceso}

El ETL se implementó en Python 3 con \texttt{pymysql} y \texttt{psycopg2},
dentro del contenedor \texttt{coop\_horizonte\_migrador}, construido sobre la
imagen \texttt{postgres:16}. El proceso abre una transacción de lectura en
MySQL con \texttt{REPEATABLE READ} y \texttt{START TRANSACTION WITH CONSISTENT
SNAPSHOT}, de modo que todas las tablas se leen desde el mismo punto en el
tiempo. Luego trunca las tablas destino, carga en orden jerárquico, reajusta
las secuencias y compara conteos. La Figura~\ref{fig:flujo-etl} resume el
flujo.

\figuraAPA{fig:flujo-etl}{Flujo del ETL y de la verificación}{%
  \begin{tikzpicture}[
      paso/.style={draw, rounded corners, minimum width=2.7cm,
                   minimum height=0.9cm, align=center, font=\scriptsize,
                   fill=green!4},
      linea/.style={-Stealth, thick}]
    \node[paso] (ext) {Extraer\\snapshot};
    \node[paso, right=0.7cm of ext] (tr) {Transformar\\tipos};
    \node[paso, right=0.7cm of tr] (load) {Cargar\\destino};
    \node[paso, right=0.7cm of load] (seq) {Reajustar\\secuencias};
    \node[paso, right=0.7cm of seq] (ver) {Verificar\\conteos};
    \draw[linea] (ext) -- (tr);
    \draw[linea] (tr) -- (load);
    \draw[linea] (load) -- (seq);
    \draw[linea] (seq) -- (ver);
  \end{tikzpicture}%
}{Elaboración propia con base en \texttt{src/migration/migrate.py}.}

\subsection{Transformaciones aplicadas}

La Tabla~\ref{tab:transformaciones} documenta las transformaciones del ETL.
Cada una responde a una incompatibilidad identificada en el diagnóstico. La
conversión de booleanos usa \texttt{bool(int(valor))}; la de fechas asigna
\texttt{ZoneInfo("America/La\_Paz")}; la de identificadores conserva el valor
original y lo inserta con \texttt{OVERRIDING SYSTEM VALUE}. Los valores
\texttt{NULL} se preservan en todas las columnas.

\begin{table}[H]
  \centering
  \caption{Transformaciones del ETL}
  \label{tab:transformaciones}
  \contenidoConFuente{%
    \begin{tablaAPA}
      \begin{tabular}{@{}L{3.0cm}L{4.4cm}L{6.8cm}@{}}
        \toprule
        \textbf{Columna origen} & \textbf{Transformación} & \textbf{Justificación} \\
        \midrule
        \texttt{estatus}, \texttt{pagado\_total} & \texttt{bool(int(valor))} & \texttt{TINYINT(1)} $\rightarrow$ \texttt{BOOLEAN}. \\
        \texttt{fecha\_*}, \texttt{fecha\_transaccion} & \texttt{valor.replace(tzinfo=America/La\_Paz)} & \texttt{DATETIME} $\rightarrow$ \texttt{TIMESTAMPTZ}. \\
        \texttt{id} & \texttt{OVERRIDING SYSTEM VALUE} & Conservar referencias y evitar colisiones. \\
        \texttt{estado}, \texttt{status}, \texttt{tipo} & Texto sin cambios; validación por \texttt{CHECK} & \texttt{ENUM} $\rightarrow$ \texttt{VARCHAR} + \texttt{CHECK}. \\
        \texttt{monto}, \texttt{saldo}, \texttt{tasa} & \texttt{Decimal} sin redondeo adicional & Mantener precisión \texttt{NUMERIC(12,2)}. \\
        \texttt{NULL} & Se conserva & No inventar valores ni aplicar defaults. \\
        \bottomrule
      \end{tabular}
    \end{tablaAPA}%
  }{Elaboración propia con base en \texttt{src/migration/migrate.py}.}
\end{table}

\subsection{Carga, idempotencia y verificación}

La carga se hace tabla por tabla, en orden de dependencias. Antes de insertar,
el ETL ejecuta \texttt{TRUNCATE} sobre las 8 tablas con \texttt{CASCADE}. No
se usa \texttt{RESTART IDENTITY} porque \texttt{admin\_dba} no es propietario
de las secuencias creadas por las identidades; en su lugar, al final se
ejecuta \texttt{setval} con el máximo ID de cada tabla. Esa decisión surgió de
un error real documentado en la bitácora
(Sección~\ref{sec:bitacora}). La carga es idempotente: se puede repetir sin
duplicar filas porque el \texttt{TRUNCATE} limpia el destino y el snapshot del
origen es consistente.

\begin{lstlisting}[language=Python,caption={Núcleo del ETL: transformación, carga y reajuste de secuencias},label={lst:etl}]
for table, columns, datetime_cols, bool_cols in TABLE_SPECS:
    mysql_cur.execute(f"SELECT {', '.join(columns)} FROM {table} ORDER BY id")
    raw_rows = mysql_cur.fetchall()
    rows = []
    for raw in raw_rows:
        rows.append(tuple(
            transform(value, col in bool_cols, col in datetime_cols)
            for col, value in zip(columns, raw)
        ))
    if rows:
        insert_sql = (
            f"INSERT INTO {table} ({', '.join(columns)}) "
            f"OVERRIDING SYSTEM VALUE VALUES %s"
        )
        execute_values(pg_cur, insert_sql, rows, page_size=500)

for table, _, _, _ in TABLE_SPECS:
    pg_cur.execute(
        f"SELECT setval(pg_get_serial_sequence(%s, 'id'), "
        f"COALESCE(MAX(id), 1), MAX(id) IS NOT NULL) FROM {table}",
        (table,),
    )
\end{lstlisting}

La verificación compara \texttt{COUNT(*)} en cada tabla del origen y del
destino. Si alguna diferencia existe, el ETL hace \texttt{rollback} y devuelve
código de salida 1. En la ejecución real, las 8 tablas coincidieron y el total
fue 1\,469 filas. La Tabla~\ref{tab:ejecucion-etl} muestra los tiempos por
tabla. La tabla más costosa fue \texttt{transacciones} (740 filas, 0.022 s),
seguida de \texttt{cuotas\_amortizacion} (516 filas, 0.014 s). El tiempo total
de carga fue inferior a 0.06 s, sin contar la conexión y la verificación.

\begin{table}[H]
  \centering
  \caption{Tiempos de carga del ETL por tabla}
  \label{tab:ejecucion-etl}
  \contenidoConFuente{%
    \begin{tablaAPA}
      \begin{tabular}{@{}L{6.4cm}R{2.4cm}R{3.0cm}@{}}
        \toprule
        \textbf{Tabla} & \textbf{Filas} & \textbf{Tiempo (s)} \\
        \midrule
        \texttt{tipos\_credito} & 5 & 0.003 \\
        \texttt{sucursales} & 5 & 0.002 \\
        \texttt{socios} & 55 & 0.003 \\
        \texttt{cuentas\_ahorro} & 66 & 0.003 \\
        \texttt{solicitudes\_credito} & 50 & 0.003 \\
        \texttt{prestamos} & 32 & 0.002 \\
        \texttt{cuotas\_amortizacion} & 516 & 0.014 \\
        \texttt{transacciones} & 740 & 0.022 \\
        \midrule
        \textbf{Total} & \textbf{1\,469} & \textbf{0.052} \\
        \bottomrule
      \end{tabular}
    \end{tablaAPA}%
  }{Elaboración propia con datos de \texttt{evidence/02\_ejecucion\_migracion.txt}.}
\end{table}

\section{Orquestación Docker y automatización}

\subsection{Servicios y aislamiento}

El archivo \texttt{docker-compose.yml} define tres servicios principales y un
perfil de documentación. La Tabla~\ref{tab:servicios} describe cada uno. Los
puertos se publican solo en \texttt{127.0.0.1}, de modo que la base no queda
expuesta a la red local. Los datos persisten en volúmenes nombrados y los
archivos de inicialización se montan como solo lectura.

\begin{table}[H]
  \centering
  \caption{Servicios de Docker Compose}
  \label{tab:servicios}
  \contenidoConFuente{%
    \begin{tablaAPA}
      \begin{tabular}{@{}L{2.8cm}L{3.0cm}L{4.2cm}L{4.4cm}@{}}
        \toprule
        \textbf{Servicio} & \textbf{Imagen} & \textbf{Función} & \textbf{Puerto / volumen} \\
        \midrule
        \texttt{db-origen} & \texttt{mysql:5.7} & Motor legacy con esquema y datos. & \texttt{127.0.0.1:3306}; \texttt{mysql\_data}. \\
        \texttt{db-destino} & \texttt{postgres:16} & Motor destino endurecido. & \texttt{127.0.0.1:5432}; \texttt{postgres\_data}. \\
        \texttt{migrador} & \texttt{postgres:16} + Python & ETL, respaldo y pruebas. & Sin puerto; red interna. \\
        \texttt{docs} & \texttt{texlive/texlive} & Compilación del informe. & Perfil \texttt{docs}. \\
        \bottomrule
      \end{tabular}
    \end{tablaAPA}%
  }{Elaboración propia con datos de \texttt{docker-compose.yml}.}
\end{table}

Los servicios comparten la red \texttt{red\_cooperativa}. El migrador resuelve
los nombres \texttt{db-origen} y \texttt{db-destino} por DNS interno de
Docker, lo que evita direcciones IP fijas. El destino monta
\texttt{postgresql.conf}, \texttt{pg\_hba.conf}, los certificados y los scripts
de inicialización. El entrypoint copia los certificados a una ruta escribible,
ajusta propietario y permisos, y delega en el entrypoint oficial de
PostgreSQL. Esa pequeña capa fue necesaria porque el volumen de datos no
permite escribir directamente en \texttt{/etc/postgresql}.

\subsection{Scripts de automatización}

La Tabla~\ref{tab:scripts} describe los cuatro scripts de soporte. Cada uno
tiene una responsabilidad única y guarda su salida en \texttt{evidence/}. El
script \texttt{lib.sh} centraliza funciones comunes: carga del \texttt{.env},
espera de servicios \emph{healthy}, ejecución de \texttt{mysql} y
\texttt{psql}, y conteos. Así se evita duplicar lógica y se facilita el
mantenimiento.

\begin{table}[H]
  \centering
  \caption{Scripts de automatización}
  \label{tab:scripts}
  \contenidoConFuente{%
    \begin{tablaAPA}
      \begin{tabular}{@{}L{4.4cm}L{9.6cm}@{}}
        \toprule
        \textbf{Script} & \textbf{Responsabilidad} \\
        \midrule
        \texttt{scripts/lib.sh} & Funciones comunes: entorno, espera, consultas y conteos. \\
        \texttt{scripts/01\_audit\_origin.sh} & Inventario, tamaño, objetos, usuarios, privilegios e incompatibilidades. \\
        \texttt{scripts/02\_migrate\_and\_backup.sh} & Respaldo del origen, pgloader + SQL, respaldo cifrado y prueba de restauración. \\
        \texttt{scripts/03\_verify\_all.sh} & 75 pruebas de hardening, roles, enmascaramiento, conteos, integridad, staging y secuencias. \\
        \bottomrule
      \end{tabular}
    \end{tablaAPA}%
  }{Elaboración propia.}
\end{table}

\subsection{Reproducibilidad}

El entorno se levanta con \texttt{docker compose up -d --build}. La primera
vez, MySQL ejecuta \texttt{01\_schema\_origin.sql} y
\texttt{02\_seed\_origin.sql}; PostgreSQL ejecuta
\texttt{01\_schema\_target.sql} y \texttt{02\_hardening\_roles.sql}. Los
healthchecks esperan a que ambos motores estén listos antes de considerar el
servicio disponible. Para repetir la migración desde cero se ejecuta
\texttt{docker compose down -v} y se vuelve a levantar; los volúmenes se
recrean y los datos se generan de nuevo con la misma semilla. Esa
reproducibilidad es la base de la confianza: cualquier integrante puede
verificar los mismos resultados sin depender de una máquina específica.
